%% fakelist.tex
%% Solving the issue of exe environments (gb4e) in inline lists.
%%
%% A fake enumerate-like environment that is just a paragraph.
%% No lists, no spacing changes, no horizontal-mode tricks.
%% exe, xlist, inexe etc. all work inside it exactly as in
%% normal paragraph text, because it IS normal paragraph text.
%%
%% Load with %% fakelist.tex
%% Solving the issue of exe environments (gb4e) in inline lists.
%%
%% A fake enumerate-like environment that is just a paragraph.
%% No lists, no spacing changes, no horizontal-mode tricks.
%% exe, xlist, inexe etc. all work inside it exactly as in
%% normal paragraph text, because it IS normal paragraph text.
%%
%% Load with %% fakelist.tex
%% Solving the issue of exe environments (gb4e) in inline lists.
%%
%% A fake enumerate-like environment that is just a paragraph.
%% No lists, no spacing changes, no horizontal-mode tricks.
%% exe, xlist, inexe etc. all work inside it exactly as in
%% normal paragraph text, because it IS normal paragraph text.
%%
%% Load with %% fakelist.tex
%% Solving the issue of exe environments (gb4e) in inline lists.
%%
%% A fake enumerate-like environment that is just a paragraph.
%% No lists, no spacing changes, no horizontal-mode tricks.
%% exe, xlist, inexe etc. all work inside it exactly as in
%% normal paragraph text, because it IS normal paragraph text.
%%
%% Load with \input{fakelist} in your preamble.
%%
%% Usage:
%%   \begin{fakenum}[(I)]
%%     \item Blah blah
%%     \begin{exe} \ex foo \end{exe}
%%     blah.
%%     \item Consider:
%%     \begin{exe}
%%       \ex This is a numbered example.
%%     \end{exe}
%%   \end{fakenum}
%%
%% Italic Labels:
%%   To italicize the label, include \itlabelfake in the optional argument:
%%   \begin{fakenum}[\itlabelfake (a)]
%%     \item This will be italicized.
%%   \end{fakenum}
%%
%% Label format (same conventions as enumerate package):
%%   (I)  uppercase roman    (i)  lowercase roman  (default)
%%   (A)  uppercase alpha    (a)  lowercase alpha
%%   (1)  arabic
%%
%% \label / \ref work on items: place \label after \item.

%% ---------------------------------------------------------------
%% Guard against double-loading
%% ---------------------------------------------------------------
\expandafter\ifx\csname fakelist@loaded\endcsname\relax
\expandafter\def\csname fakelist@loaded\endcsname{}
\else\expandafter\endinput\fi
\def\itlabelfake{}

\makeatletter

%% ---------------------------------------------------------------
%% Counter
%% ---------------------------------------------------------------

\newcounter{fakenum}

%% ---------------------------------------------------------------
%% Label format parser
%%
%% Finds the counter-style letter (A/a/I/i/1) in the spec,
%% sets \fl@fmt accordingly, and splits spec into \fl@pre/\fl@post.
%% Uses \@tfor to walk tokens — no recursion, no \if nesting issues.
%% ---------------------------------------------------------------

\def\fl@pre{}
\def\fl@fmt{\roman{fakenum}}
\def\fl@post{}
\newif\if@fl@found

\def\fl@parse#1{%
  \def\fl@pre{}%
  \def\fl@post{}%
  \def\fl@fmt{\roman{fakenum}}%
  \@fl@foundfalse
  \@tfor\fl@tok:=#1\do{%
    \if@fl@found
      \edef\fl@post{\fl@post\fl@tok}%
    \else
      \ifx\fl@tok\fl@A \def\fl@fmt{\Alph{fakenum}}\@fl@foundtrue\else
      \ifx\fl@tok\fl@a \def\fl@fmt{\alph{fakenum}}\@fl@foundtrue\else
      \ifx\fl@tok\fl@I \def\fl@fmt{\Roman{fakenum}}\@fl@foundtrue\else
      \ifx\fl@tok\fl@i \def\fl@fmt{\roman{fakenum}}\@fl@foundtrue\else
      \ifx\fl@tok\fl@one \def\fl@fmt{\arabic{fakenum}}\@fl@foundtrue\else
        \edef\fl@pre{\fl@pre\fl@tok}%
      \fi\fi\fi\fi\fi
    \fi
  }%
}

\def\fl@A{A}\def\fl@a{a}\def\fl@I{I}\def\fl@i{i}\def\fl@one{1}

%% ---------------------------------------------------------------
%% \item replacement
%% ---------------------------------------------------------------

\def\fakenumfont#1{#1}
\def\fl@item{%
  \ifnum\@listdepth>\z@
    % Inside a real list (e.g. exe): use the real \item
    \fl@saveditem
  \else
    % Top level in fakenum: print our fake label
    \refstepcounter{fakenum}%
    \protected@edef\@currentlabel{\fl@pre\fl@fmt\fl@post}%
    \fl@pre\fakenumfont{\fl@fmt}\fl@post~%\enskip%\hspace{.5em}%
  \fi
}
\def\fl@item{%
  \ifnum\@listdepth>\z@
    \fl@saveditem
  \else
    \refstepcounter{fakenum}%
    \protected@edef\@currentlabel{\fl@pre\fl@fmt\fl@post}%
    % This line ensures the font command wraps the formatted counter
    \fl@pre\fakenumfont{\fl@fmt}\fl@post~% 
  \fi
}
\def\fl@item{%
  \ifnum\@listdepth>\z@
    \fl@saveditem
  \else
    \refstepcounter{fakenum}%
    \protected@edef\@currentlabel{\fl@pre\fl@fmt\fl@post}%
    % This line ensures the font command wraps the formatted counter
    \fl@pre\fakenumfont{\fl@fmt}\fl@post~% 
  \fi
}
%% ---------------------------------------------------------------
%% fakenum environment
%% ---------------------------------------------------------------

% \newenvironment{fakenum}[1][(i)]{%
%   \setcounter{fakenum}{0}%
%   \fl@parse{#1}%
%   \let\fl@saveditem\item
%   \let\item\fl@item
%   \ignorespaces
% }{%
%   \let\item\fl@saveditem
% }
% \renewenvironment{fakenum}[1][(i)]{%
%   \setcounter{fakenum}{0}%
%   \def\fakenumfont##1{##1}% Reset font to default for each environment
%   #1% Execute commands like \itlabelfake if present in the argument
%   \fl@parse{#1}% Parse the string for labels
%   \let\fl@saveditem\item
%   \let\item\fl@item
%   \ignorespaces
% }{%
%   \let\item\fl@item
% }
% \renewenvironment{fakenum}[1][(i)]{%
%   \setcounter{fakenum}{0}%
%   \def\fakenumfont##1{##1}% Reset to default (upright) 
%   #1% Execute \itlabelfake if it's in the brackets
%   \fl@parse{#1}% Run the parser [cite: 9]
%   \let\fl@saveditem\item
%   \let\item\fl@item
%   \ignorespaces
% }{%
%   \let\item\fl@saveditem
% }

\newenvironment{fakenum}[1][(i)]{%
  \setcounter{fakenum}{0}%
  % 1. Default: Upright font
  \def\fakenumfont##1{##1}% 
  % 2. Check if \itlabelfake was typed in the brackets
  \tokenize@check{#1}{\itlabelfake}{\def\fakenumfont##1{(\textit{##1})}}%
  % 3. Run the original parser 
  \fl@parse{#1}% 
  \let\fl@saveditem\item
  \let\item\fl@item
  \ignorespaces
}{%
  \let\item\fl@saveditem
}

% Helper to check for the command without expanding it
\def\tokenize@check#1#2#3{%
  \def\temp@a{#1}\def\temp@b{#2}%
  \ifx\temp@a\temp@b #3\fi % Simple check if only \itlabelfake is present
  % Or more simply, just check if the argument contains it:
  \inner@check#1\itlabelfake\@@@{#3}%
}
\def\inner@check#1\itlabelfake#2\@@@#3{\ifx\relax#2\relax\else#3\fi}

\makeatother

\endinput
 in your preamble.
%%
%% Usage:
%%   \begin{fakenum}[(I)]
%%     \item Blah blah
%%     \begin{exe} \ex foo \end{exe}
%%     blah.
%%     \item Consider:
%%     \begin{exe}
%%       \ex This is a numbered example.
%%     \end{exe}
%%   \end{fakenum}
%%
%% Italic Labels:
%%   To italicize the label, include \itlabelfake in the optional argument:
%%   \begin{fakenum}[\itlabelfake (a)]
%%     \item This will be italicized.
%%   \end{fakenum}
%%
%% Label format (same conventions as enumerate package):
%%   (I)  uppercase roman    (i)  lowercase roman  (default)
%%   (A)  uppercase alpha    (a)  lowercase alpha
%%   (1)  arabic
%%
%% \label / \ref work on items: place \label after \item.

%% ---------------------------------------------------------------
%% Guard against double-loading
%% ---------------------------------------------------------------
\expandafter\ifx\csname fakelist@loaded\endcsname\relax
\expandafter\def\csname fakelist@loaded\endcsname{}
\else\expandafter\endinput\fi
\def\itlabelfake{}

\makeatletter

%% ---------------------------------------------------------------
%% Counter
%% ---------------------------------------------------------------

\newcounter{fakenum}

%% ---------------------------------------------------------------
%% Label format parser
%%
%% Finds the counter-style letter (A/a/I/i/1) in the spec,
%% sets \fl@fmt accordingly, and splits spec into \fl@pre/\fl@post.
%% Uses \@tfor to walk tokens — no recursion, no \if nesting issues.
%% ---------------------------------------------------------------

\def\fl@pre{}
\def\fl@fmt{\roman{fakenum}}
\def\fl@post{}
\newif\if@fl@found

\def\fl@parse#1{%
  \def\fl@pre{}%
  \def\fl@post{}%
  \def\fl@fmt{\roman{fakenum}}%
  \@fl@foundfalse
  \@tfor\fl@tok:=#1\do{%
    \if@fl@found
      \edef\fl@post{\fl@post\fl@tok}%
    \else
      \ifx\fl@tok\fl@A \def\fl@fmt{\Alph{fakenum}}\@fl@foundtrue\else
      \ifx\fl@tok\fl@a \def\fl@fmt{\alph{fakenum}}\@fl@foundtrue\else
      \ifx\fl@tok\fl@I \def\fl@fmt{\Roman{fakenum}}\@fl@foundtrue\else
      \ifx\fl@tok\fl@i \def\fl@fmt{\roman{fakenum}}\@fl@foundtrue\else
      \ifx\fl@tok\fl@one \def\fl@fmt{\arabic{fakenum}}\@fl@foundtrue\else
        \edef\fl@pre{\fl@pre\fl@tok}%
      \fi\fi\fi\fi\fi
    \fi
  }%
}

\def\fl@A{A}\def\fl@a{a}\def\fl@I{I}\def\fl@i{i}\def\fl@one{1}

%% ---------------------------------------------------------------
%% \item replacement
%% ---------------------------------------------------------------

\def\fakenumfont#1{#1}
\def\fl@item{%
  \ifnum\@listdepth>\z@
    % Inside a real list (e.g. exe): use the real \item
    \fl@saveditem
  \else
    % Top level in fakenum: print our fake label
    \refstepcounter{fakenum}%
    \protected@edef\@currentlabel{\fl@pre\fl@fmt\fl@post}%
    \fl@pre\fakenumfont{\fl@fmt}\fl@post~%\enskip%\hspace{.5em}%
  \fi
}
\def\fl@item{%
  \ifnum\@listdepth>\z@
    \fl@saveditem
  \else
    \refstepcounter{fakenum}%
    \protected@edef\@currentlabel{\fl@pre\fl@fmt\fl@post}%
    % This line ensures the font command wraps the formatted counter
    \fl@pre\fakenumfont{\fl@fmt}\fl@post~% 
  \fi
}
\def\fl@item{%
  \ifnum\@listdepth>\z@
    \fl@saveditem
  \else
    \refstepcounter{fakenum}%
    \protected@edef\@currentlabel{\fl@pre\fl@fmt\fl@post}%
    % This line ensures the font command wraps the formatted counter
    \fl@pre\fakenumfont{\fl@fmt}\fl@post~% 
  \fi
}
%% ---------------------------------------------------------------
%% fakenum environment
%% ---------------------------------------------------------------

% \newenvironment{fakenum}[1][(i)]{%
%   \setcounter{fakenum}{0}%
%   \fl@parse{#1}%
%   \let\fl@saveditem\item
%   \let\item\fl@item
%   \ignorespaces
% }{%
%   \let\item\fl@saveditem
% }
% \renewenvironment{fakenum}[1][(i)]{%
%   \setcounter{fakenum}{0}%
%   \def\fakenumfont##1{##1}% Reset font to default for each environment
%   #1% Execute commands like \itlabelfake if present in the argument
%   \fl@parse{#1}% Parse the string for labels
%   \let\fl@saveditem\item
%   \let\item\fl@item
%   \ignorespaces
% }{%
%   \let\item\fl@item
% }
% \renewenvironment{fakenum}[1][(i)]{%
%   \setcounter{fakenum}{0}%
%   \def\fakenumfont##1{##1}% Reset to default (upright) 
%   #1% Execute \itlabelfake if it's in the brackets
%   \fl@parse{#1}% Run the parser [cite: 9]
%   \let\fl@saveditem\item
%   \let\item\fl@item
%   \ignorespaces
% }{%
%   \let\item\fl@saveditem
% }

\newenvironment{fakenum}[1][(i)]{%
  \setcounter{fakenum}{0}%
  % 1. Default: Upright font
  \def\fakenumfont##1{##1}% 
  % 2. Check if \itlabelfake was typed in the brackets
  \tokenize@check{#1}{\itlabelfake}{\def\fakenumfont##1{(\textit{##1})}}%
  % 3. Run the original parser 
  \fl@parse{#1}% 
  \let\fl@saveditem\item
  \let\item\fl@item
  \ignorespaces
}{%
  \let\item\fl@saveditem
}

% Helper to check for the command without expanding it
\def\tokenize@check#1#2#3{%
  \def\temp@a{#1}\def\temp@b{#2}%
  \ifx\temp@a\temp@b #3\fi % Simple check if only \itlabelfake is present
  % Or more simply, just check if the argument contains it:
  \inner@check#1\itlabelfake\@@@{#3}%
}
\def\inner@check#1\itlabelfake#2\@@@#3{\ifx\relax#2\relax\else#3\fi}

\makeatother

\endinput
 in your preamble.
%%
%% Usage:
%%   \begin{fakenum}[(I)]
%%     \item Blah blah
%%     \begin{exe} \ex foo \end{exe}
%%     blah.
%%     \item Consider:
%%     \begin{exe}
%%       \ex This is a numbered example.
%%     \end{exe}
%%   \end{fakenum}
%%
%% Italic Labels:
%%   To italicize the label, include \itlabelfake in the optional argument:
%%   \begin{fakenum}[\itlabelfake (a)]
%%     \item This will be italicized.
%%   \end{fakenum}
%%
%% Label format (same conventions as enumerate package):
%%   (I)  uppercase roman    (i)  lowercase roman  (default)
%%   (A)  uppercase alpha    (a)  lowercase alpha
%%   (1)  arabic
%%
%% \label / \ref work on items: place \label after \item.

%% ---------------------------------------------------------------
%% Guard against double-loading
%% ---------------------------------------------------------------
\expandafter\ifx\csname fakelist@loaded\endcsname\relax
\expandafter\def\csname fakelist@loaded\endcsname{}
\else\expandafter\endinput\fi
\def\itlabelfake{}

\makeatletter

%% ---------------------------------------------------------------
%% Counter
%% ---------------------------------------------------------------

\newcounter{fakenum}

%% ---------------------------------------------------------------
%% Label format parser
%%
%% Finds the counter-style letter (A/a/I/i/1) in the spec,
%% sets \fl@fmt accordingly, and splits spec into \fl@pre/\fl@post.
%% Uses \@tfor to walk tokens — no recursion, no \if nesting issues.
%% ---------------------------------------------------------------

\def\fl@pre{}
\def\fl@fmt{\roman{fakenum}}
\def\fl@post{}
\newif\if@fl@found

\def\fl@parse#1{%
  \def\fl@pre{}%
  \def\fl@post{}%
  \def\fl@fmt{\roman{fakenum}}%
  \@fl@foundfalse
  \@tfor\fl@tok:=#1\do{%
    \if@fl@found
      \edef\fl@post{\fl@post\fl@tok}%
    \else
      \ifx\fl@tok\fl@A \def\fl@fmt{\Alph{fakenum}}\@fl@foundtrue\else
      \ifx\fl@tok\fl@a \def\fl@fmt{\alph{fakenum}}\@fl@foundtrue\else
      \ifx\fl@tok\fl@I \def\fl@fmt{\Roman{fakenum}}\@fl@foundtrue\else
      \ifx\fl@tok\fl@i \def\fl@fmt{\roman{fakenum}}\@fl@foundtrue\else
      \ifx\fl@tok\fl@one \def\fl@fmt{\arabic{fakenum}}\@fl@foundtrue\else
        \edef\fl@pre{\fl@pre\fl@tok}%
      \fi\fi\fi\fi\fi
    \fi
  }%
}

\def\fl@A{A}\def\fl@a{a}\def\fl@I{I}\def\fl@i{i}\def\fl@one{1}

%% ---------------------------------------------------------------
%% \item replacement
%% ---------------------------------------------------------------

\def\fakenumfont#1{#1}
\def\fl@item{%
  \ifnum\@listdepth>\z@
    % Inside a real list (e.g. exe): use the real \item
    \fl@saveditem
  \else
    % Top level in fakenum: print our fake label
    \refstepcounter{fakenum}%
    \protected@edef\@currentlabel{\fl@pre\fl@fmt\fl@post}%
    \fl@pre\fakenumfont{\fl@fmt}\fl@post~%\enskip%\hspace{.5em}%
  \fi
}
\def\fl@item{%
  \ifnum\@listdepth>\z@
    \fl@saveditem
  \else
    \refstepcounter{fakenum}%
    \protected@edef\@currentlabel{\fl@pre\fl@fmt\fl@post}%
    % This line ensures the font command wraps the formatted counter
    \fl@pre\fakenumfont{\fl@fmt}\fl@post~% 
  \fi
}
\def\fl@item{%
  \ifnum\@listdepth>\z@
    \fl@saveditem
  \else
    \refstepcounter{fakenum}%
    \protected@edef\@currentlabel{\fl@pre\fl@fmt\fl@post}%
    % This line ensures the font command wraps the formatted counter
    \fl@pre\fakenumfont{\fl@fmt}\fl@post~% 
  \fi
}
%% ---------------------------------------------------------------
%% fakenum environment
%% ---------------------------------------------------------------

% \newenvironment{fakenum}[1][(i)]{%
%   \setcounter{fakenum}{0}%
%   \fl@parse{#1}%
%   \let\fl@saveditem\item
%   \let\item\fl@item
%   \ignorespaces
% }{%
%   \let\item\fl@saveditem
% }
% \renewenvironment{fakenum}[1][(i)]{%
%   \setcounter{fakenum}{0}%
%   \def\fakenumfont##1{##1}% Reset font to default for each environment
%   #1% Execute commands like \itlabelfake if present in the argument
%   \fl@parse{#1}% Parse the string for labels
%   \let\fl@saveditem\item
%   \let\item\fl@item
%   \ignorespaces
% }{%
%   \let\item\fl@item
% }
% \renewenvironment{fakenum}[1][(i)]{%
%   \setcounter{fakenum}{0}%
%   \def\fakenumfont##1{##1}% Reset to default (upright) 
%   #1% Execute \itlabelfake if it's in the brackets
%   \fl@parse{#1}% Run the parser [cite: 9]
%   \let\fl@saveditem\item
%   \let\item\fl@item
%   \ignorespaces
% }{%
%   \let\item\fl@saveditem
% }

\newenvironment{fakenum}[1][(i)]{%
  \setcounter{fakenum}{0}%
  % 1. Default: Upright font
  \def\fakenumfont##1{##1}% 
  % 2. Check if \itlabelfake was typed in the brackets
  \tokenize@check{#1}{\itlabelfake}{\def\fakenumfont##1{(\textit{##1})}}%
  % 3. Run the original parser 
  \fl@parse{#1}% 
  \let\fl@saveditem\item
  \let\item\fl@item
  \ignorespaces
}{%
  \let\item\fl@saveditem
}

% Helper to check for the command without expanding it
\def\tokenize@check#1#2#3{%
  \def\temp@a{#1}\def\temp@b{#2}%
  \ifx\temp@a\temp@b #3\fi % Simple check if only \itlabelfake is present
  % Or more simply, just check if the argument contains it:
  \inner@check#1\itlabelfake\@@@{#3}%
}
\def\inner@check#1\itlabelfake#2\@@@#3{\ifx\relax#2\relax\else#3\fi}

\makeatother

\endinput
 in your preamble.
%%
%% Usage:
%%   \begin{fakenum}[(I)]
%%     \item Blah blah
%%     \begin{exe} \ex foo \end{exe}
%%     blah.
%%     \item Consider:
%%     \begin{exe}
%%       \ex This is a numbered example.
%%     \end{exe}
%%   \end{fakenum}
%%
%% Italic Labels:
%%   To italicize the label, include \itlabelfake in the optional argument:
%%   \begin{fakenum}[\itlabelfake (a)]
%%     \item This will be italicized.
%%   \end{fakenum}
%%
%% Label format (same conventions as enumerate package):
%%   (I)  uppercase roman    (i)  lowercase roman  (default)
%%   (A)  uppercase alpha    (a)  lowercase alpha
%%   (1)  arabic
%%
%% \label / \ref work on items: place \label after \item.

%% ---------------------------------------------------------------
%% Guard against double-loading
%% ---------------------------------------------------------------
\expandafter\ifx\csname fakelist@loaded\endcsname\relax
\expandafter\def\csname fakelist@loaded\endcsname{}
\else\expandafter\endinput\fi
\def\itlabelfake{}

\makeatletter

%% ---------------------------------------------------------------
%% Counter
%% ---------------------------------------------------------------

\newcounter{fakenum}

%% ---------------------------------------------------------------
%% Label format parser
%%
%% Finds the counter-style letter (A/a/I/i/1) in the spec,
%% sets \fl@fmt accordingly, and splits spec into \fl@pre/\fl@post.
%% Uses \@tfor to walk tokens — no recursion, no \if nesting issues.
%% ---------------------------------------------------------------

\def\fl@pre{}
\def\fl@fmt{\roman{fakenum}}
\def\fl@post{}
\newif\if@fl@found

\def\fl@parse#1{%
  \def\fl@pre{}%
  \def\fl@post{}%
  \def\fl@fmt{\roman{fakenum}}%
  \@fl@foundfalse
  \@tfor\fl@tok:=#1\do{%
    \if@fl@found
      \edef\fl@post{\fl@post\fl@tok}%
    \else
      \ifx\fl@tok\fl@A \def\fl@fmt{\Alph{fakenum}}\@fl@foundtrue\else
      \ifx\fl@tok\fl@a \def\fl@fmt{\alph{fakenum}}\@fl@foundtrue\else
      \ifx\fl@tok\fl@I \def\fl@fmt{\Roman{fakenum}}\@fl@foundtrue\else
      \ifx\fl@tok\fl@i \def\fl@fmt{\roman{fakenum}}\@fl@foundtrue\else
      \ifx\fl@tok\fl@one \def\fl@fmt{\arabic{fakenum}}\@fl@foundtrue\else
        \edef\fl@pre{\fl@pre\fl@tok}%
      \fi\fi\fi\fi\fi
    \fi
  }%
}

\def\fl@A{A}\def\fl@a{a}\def\fl@I{I}\def\fl@i{i}\def\fl@one{1}

%% ---------------------------------------------------------------
%% \item replacement
%% ---------------------------------------------------------------

\def\fakenumfont#1{#1}
\def\fl@item{%
  \ifnum\@listdepth>\z@
    % Inside a real list (e.g. exe): use the real \item
    \fl@saveditem
  \else
    % Top level in fakenum: print our fake label
    \refstepcounter{fakenum}%
    \protected@edef\@currentlabel{\fl@pre\fl@fmt\fl@post}%
    \fl@pre\fakenumfont{\fl@fmt}\fl@post~%\enskip%\hspace{.5em}%
  \fi
}
\def\fl@item{%
  \ifnum\@listdepth>\z@
    \fl@saveditem
  \else
    \refstepcounter{fakenum}%
    \protected@edef\@currentlabel{\fl@pre\fl@fmt\fl@post}%
    % This line ensures the font command wraps the formatted counter
    \fl@pre\fakenumfont{\fl@fmt}\fl@post~% 
  \fi
}
\def\fl@item{%
  \ifnum\@listdepth>\z@
    \fl@saveditem
  \else
    \refstepcounter{fakenum}%
    \protected@edef\@currentlabel{\fl@pre\fl@fmt\fl@post}%
    % This line ensures the font command wraps the formatted counter
    \fl@pre\fakenumfont{\fl@fmt}\fl@post~% 
  \fi
}
%% ---------------------------------------------------------------
%% fakenum environment
%% ---------------------------------------------------------------

% \newenvironment{fakenum}[1][(i)]{%
%   \setcounter{fakenum}{0}%
%   \fl@parse{#1}%
%   \let\fl@saveditem\item
%   \let\item\fl@item
%   \ignorespaces
% }{%
%   \let\item\fl@saveditem
% }
% \renewenvironment{fakenum}[1][(i)]{%
%   \setcounter{fakenum}{0}%
%   \def\fakenumfont##1{##1}% Reset font to default for each environment
%   #1% Execute commands like \itlabelfake if present in the argument
%   \fl@parse{#1}% Parse the string for labels
%   \let\fl@saveditem\item
%   \let\item\fl@item
%   \ignorespaces
% }{%
%   \let\item\fl@item
% }
% \renewenvironment{fakenum}[1][(i)]{%
%   \setcounter{fakenum}{0}%
%   \def\fakenumfont##1{##1}% Reset to default (upright) 
%   #1% Execute \itlabelfake if it's in the brackets
%   \fl@parse{#1}% Run the parser [cite: 9]
%   \let\fl@saveditem\item
%   \let\item\fl@item
%   \ignorespaces
% }{%
%   \let\item\fl@saveditem
% }

\newenvironment{fakenum}[1][(i)]{%
  \setcounter{fakenum}{0}%
  % 1. Default: Upright font
  \def\fakenumfont##1{##1}% 
  % 2. Check if \itlabelfake was typed in the brackets
  \tokenize@check{#1}{\itlabelfake}{\def\fakenumfont##1{(\textit{##1})}}%
  % 3. Run the original parser 
  \fl@parse{#1}% 
  \let\fl@saveditem\item
  \let\item\fl@item
  \ignorespaces
}{%
  \let\item\fl@saveditem
}

% Helper to check for the command without expanding it
\def\tokenize@check#1#2#3{%
  \def\temp@a{#1}\def\temp@b{#2}%
  \ifx\temp@a\temp@b #3\fi % Simple check if only \itlabelfake is present
  % Or more simply, just check if the argument contains it:
  \inner@check#1\itlabelfake\@@@{#3}%
}
\def\inner@check#1\itlabelfake#2\@@@#3{\ifx\relax#2\relax\else#3\fi}

\makeatother

\endinput
